\section{Thoughts on Initiation}

\begin{quotex}
Noli foras ire, in te ipsum redi: in interiore hominis habitat veritas. 
\flright{St. Augustine}
\end{quotex}

\textbf{Rene Guenon} has written extensively on the conditions and facts of initiation, but hardly enough on what occurs after initiation. I offer, here, some thoughts on initiation, and bring up the question of its necessity and sufficiency. Unless a man prefers to be tossed around in life like Brownian motion, the final cause is always the most important. Hence, the first topic will revolve around an understanding of the goal of initiation, followed by the requirements and availability of initiation, and finally some more practical and less theoretical issues.

\paragraph{The Goal of Initiation}


I will discuss this in terms of the dominant European spiritual tradition since most readers will be familiar with that perspective. Briefly stated, a man needs to follow certain moral norms, participate in religious rites, and adhere to a creed. Then, at death, he will be judged, and ``sent'' to a place that, as a matter of justice, is his proper punishment or reward. The highest reward is the vision of God in Heaven. This is the \emph{exoteric} view, since everything is expressed and understood extrinsically, analogously to human and material affairs.

From the \emph{esoteric} point of view, the same teaching is understood symbolically. The exoteric mind can only ``picture'' the teachings, i.e., he ``sees'' them as similar to the physical world of persons and things, although operating in a ``spiritual'' world. The esoteric mind, on the other hand, also ``sees'' these teachings, but he sees interiorly. This is quite different from the philosopher or theologian who want to conceptualize the teachings and deal with them rationally.

Esoterically, Heaven, Hell, and Purgatory are states of being, that are experienced now, not \emph{post mortem}. The path is to travel through these states, from the lowest to the highest, in this life. The vision of God, or the Supreme Identity, is the highest point on that path, when all possibilities have been realized and all worlds transcended. Moral norms are the means to become aware of the states of one's being; the resistance to them reveals something about oneself. The state a man is in is not imposed upon him by a judge, but it is where he belongs by justice, since his state is the result of his own efforts.

This, then, is the goal of initiation; while it provides the initial spiritual impulse, the ongoing support of the organization is even more necessary. The path is arduous; the initiate will discover unpleasant facts about himself; boredom, fear, and anxiety will tempt him to falter. But, this is the greater battle that initiates need to fight.

\paragraph{The Conditions of Initiation}


Guenon points out that a valid initiation requires

\begin{enumerate}


\item An unbroken chain of initiation whose originary source was supernatural,

\item An attachment to one of the true Traditions.
\end{enumerate}


These requirements are difficult to meet currently in the West, at least in regards to an authentically Western tradition, although various Oriental groups have set up bases in Europe and North America over the last couple of generations. Pre-schism Christianity satisfied both conditions and there were clearly initiatic organizations in the Middle Ages, and most likely, remnants still existing today.

In any case, it is important to point out that initiation is not simply open to anyone. There is an extended period of preparation necessary, even if it is primarily and intellectual preparation. The spiritual literature abounds with the difficulties that seekers experience when looking for an initiation. For example, in the Pythagorean group, the seeker had to be silent for several years; Sufi groups may require 1001 days of menial tasks prior to initiation. In my opinion, Guenon fails to emphasize this enough.

Furthermore, applicant himself has to satisfy certain conditions and be free of certain physical defects. These may seem arbitrary and unjust to today's egalitarian spirit, but there are valid reasons for them, as Guenon discusses in Perspectives on Initiation. There is no point to mention them all here, especially since they may vary from group to group, but there is one that may surprise readers.

Priestly ordination is an example Guenon uses, even if today it no longer seems to qualify as a valid initiatory path. St Thomas Aquinas pointed out that effeminacy in speech, mannerisms, or posture would disqualify a man from becoming a priest. This requirement would offend most Christians today although it is certainly consistent with the Traditional ideal of a virile priesthood. Nevertheless, the fruits of a feminized church are too obvious to mention.

\paragraph{Practical Issues}


First of all, there is the question of whether initiation is sufficient to lead a man along the true spiritual path. Well, Guenon provided us the example of Inayat Khan, who took initiations with multiple orthodox Sufi groups. Yet, he still deviated from orthodox tradition when he formed his own group. It eventually became associated with the Human Potential Movement in North America and his grandson is the current leader of the group Mr. Khan founded. Why it is unorthodox is another topic.

If initiation is not sufficient, is it necessary? This is more subtle than it appears. It is certainly necessary to understand the goal of the initiatory path, the means to accomplish it, and the metaphysical teachings underlying it. Yet the chain is both human and superhuman. As such, initiation may be made possible through superhuman or preternatural means. For example, Tomberg claimed to be guided by long deceased Hermetists. This is outrageous only If one regards this world as a closed system.

It must be pointed out that any teachings from such an initiation applies only to the person receiving it and is not binding on anyone else. There are too many cases to mention of men or women who receive a minor revelation and through ego inflation go on to found new religions. Ultimately, orthodox traditional teachings must be the one and only touchstone.

On a more practical level, Guenon does not make clear the distinction between a valid and a licit initiation. A valid initiation is performed by a person with the power to confer it, but a licit initiation is performed by a person with the authority to confer it. Ideally, an initiation is both valid and licit, that is, it falls under the jurisdiction of a traditional spiritual form. Nevertheless, there are certainly cases of groups under the leadership of someone with a valid initiation, but illicit since it operates independently of a legitimate spiritual authority. Some of the figures that intrigued Julius Evola\footnote{\url{https://gornahoor.net/?p=4628}} may have been of this nature.

Although there are dangers in such groups, this is mitigated is one understands the goal of initiation and true metaphysics. The dangers may, in fact, be offset by the opportunities gained to make the teachings operative rather than only speculative. Given the concrete conditions a man seeking spiritual gnosis finds himself today, that may be the best, or only, possibility for him. If he starts from Smalltown, USA, a valid and licit initiation may be like a first class plane ticket to the bright lights of Hollywood. Yet, as we pointed out, and will do so again, there is no guarantee the plane will arrive safely.

Some of us may have to be content with a blue-collar initiation, more like a bus ticket to St. Louis. In any event, it is better than staying home.

\flrightit{Posted on 2012-08-15 by Cologero}

\begin{center}* * *\end{center}

\begin{footnotesize}
\begin{sffamily}
\texttt{Ferdinand Gombault on 2012-08-16 at 03:23 said:}

Highly interesting interpretation of the concept of initiation. It is actually quite unclear, what ``traditionalist'' authors understand under initiation and spiritual realiziation. In this context, Evolas analysis of the ancient doctrine of the Paths in the Afterlife. It shows that for Evola (or for Tradition in his understanding) the afterlife was not a mere symbol for states of the being during the time the corporal body is alive. In particular, there exists a relationship between initiation and the ``path of the gods''.

In ``Sophia: The Journal of Traditional Studies'' appeared an article titled ``René Guénon and the question of initiation'' by Timothy Scott, downloadable on Mr. Scotts website:

\textit{http://timothyscott.com.au/scholarship.html}

under articles, a quite crtitical discussion of Guenons teachings.

\hfill

\texttt{Graham on 2012-08-16 at 11:20 said:}

Ferdinand,

Regarding this comment:

``It shows that for Evola (or for Tradition in his understanding) the afterlife was not a mere symbol for states of the being during the time the corporal body is alive.''

I don't believe anyone has said that the afterlife, in the exoteric sense, is a `mere symbol'. That the afterlife -- purgatory and the heavens -- can be traversed post mortem, and that it can be traversed here and now, are perfectly compatible doctrines. Catholic mystical theology speaks of the purgative, illuminative, and unitive states which can be attained in the present life. This idea is even an element of Christian exoterism: ``Know you not that all we, who are baptized in Christ Jesus, are baptized in his death? For we are buried together with him by baptism into death; that as Christ is risen from the dead by the glory of the Father, so we also may walk in newness of life.'' If Christians are baptised into death, then they commence the afterlife while on earth.

Cologero,

The assertion that one's state is ``the result of one's own efforts'' needs to be clarified. Guenon was I think too cavalier regarding this doctrine, of the purely active effort of initiation. Schuon even takes him to task for this in his `Observations'. If we are serious about orthodoxy, and I am, this assertion borders on Pelagianism. I'd appreciate your perspective on this.

\hfill

\texttt{Michael on 2012-08-16 at 13:39 said:}

``These requirements are difficult to meet currently in the West.'' Presumably the medieval Church was a valid channel of initiation, but why does the present-day church fail the test?

\hfill

\texttt{Graham on 2012-08-16 at 14:46 said:}

Michael,

Good question. It seems to be a case of ``by their fruits you shall know them'' -- and the present-day visible Church does, on some level, fail the test. Why? Schuon, in his Observations, actually mocks Guenon's idea that God might lessen or withdraw the efficacy of sacraments He established, speaking particularly of Baptism. In lieu of Guenon's hypothesis, Schuon suggests -- if I understand him -- the quality of the people baptized and the loss of actualizing techniques as explanations more in line with principles.

\hfill

\texttt{apeiron on 2012-08-16 at 17:17 said:}

To a greater or lesser extent, initiation has now been achieved by the operator who is liberated in the full light of their own self-realized godhood. Now what?

Is it understood by the Author that post-mortem consciousness is possible in the afterlife or is this ``exoteric'' understanding of altered states of being just some symbolism we experience through life lived and thus eternal unconsciousness awaits the physical death?

\hfill

\texttt{Cologero on 2012-08-16 at 19:52 said:}

Certainly interesting comments to which I can only offer an opinion.

RE: whether post-mortem consciousness is possible in the afterlife

According to Guenon, the body is one modality of the human state. Hence, post-mortem the other modalities may persist indefinitely and they are the conscious part.

The symbolism is the expression of the modalities and states of being in physical terms, not that symbolism itself nullifies the reality it is intended to express.

The difference is that the being is passive in that state and cannot actively progress.

\hfill

\texttt{Cologero on 2012-08-16 at 19:59 said:}

Thank you for the link, Mr. Gombault.

What ``traditionalist'' authors understand by spiritual realiziation doesn't seem unclear. Even Mr. Scott writes: ``The final goal of the initiatic work is that the initiate transcend individuality in achieving Deliverance, which is the state of Supreme Identity with the Reality.''

This is what I tried to say.

Apparently I did not express myself clearly, since you bring up the same issue as apeiron. The afterlife itself is not merely symbolic. The symbols describe actual states of consciousness. The difference is whether the states are traversed passively and only after death, or whether there is a path to follow in an active way, so those states can be experienced during physical life.

\hfill

\texttt{Cologero on 2012-08-16 at 20:04 said:}

Graham, I think I'll leave it up to you to find the best possible interpretation that evades the charge of Pelagianism. Some helpful hints:

I referred to initiation rites (or sacraments), as well as following the exoteric tradition. The issue of why some men decide to follow such a path and why even fewer may be qualified bring up the complex notions of predestination and predilection respectively.

Dante relied on guides and was chosen by the intercession of Beatrice, yet he had to make the trip himself.

\hfill

\texttt{Ferdinand Gombault on 2012-08-17 at 04:38 said:}

Dear Cologero,

thank your for this clarifications. I have to correct myself. I did not want to insinuate that you are considering afterlife as merely symbolic; this was a hint to modern understandig of various traditional elements in general,

What I was trying to say is that for me after years of reading various authors in this field I still am not satisfied by the illustrations in this regard. The responsibilty for this is clearly on my side.

This deficiency in particular applies to the ``afterlife'', which I mentioned. You are right, there are clear explanations not only by ``traditionalist'' but also from ``classic'' authors, consistent with each other to a large extent. (By no means this differentiation is meant pejoratively, but rather as a chronological distinction.) It seems obvious that the state of Supreme Identity with the Reality equals ``immortality'' since Reality is eternal and, since that this Supreme Identity profoundly differs from ordinary self-awareness, cannot be described in terms of common conditional experience. Therefore, post-mortem consciousness in the afterlife is thoroughly different from consciousness during life. For present men, regrettably living in the ruins of modernity, all this is hard to figure out.

\hfill

\end{sffamily}
\end{footnotesize}

